\section{Negative critical half-integers}\label{sec:threshold}

It remains to treat
\begin{equation}\label{eq:threshold-parameter}
 \delta=-m-\frac12,
 \qquad m=1,2,\ldots.
\end{equation}
Then
\[
 \alpha=1+\delta=\frac12-m,
 \qquad
 \sin(\pi\alpha)=(-1)^m.
\]
It is convenient in this section to write
\begin{equation}\label{eq:KTa}
 K_a:=\frac1\pi\Hh_a,
 \qquad
 T_a:=\id-K_a,
 \qquad
 T_{N,a}:=P_NT_aP_N\big|_{P_N\ell^2(\Nzero)}.
\end{equation}
Thus $A_\alpha=(-1)^mK_\alpha$.  The determinant whose absolutely continuous
spectrum reaches zero is $D_N^-$ if $m$ is even and $D_N^+$ if $m$ is odd; in
either case it equals $\det T_{N,\alpha}$.

The proof requires a weak quantitative form of endpoint coercivity.  For a
sequence $0<\gamma_N\le1$, we write $\gamma_N=N^{-o(1)}$ when
\begin{equation}\label{eq:subpolynomial-definition}
 -\log\gamma_N=o(\log N),
\end{equation}
equivalently, $\gamma_N^{-1}=N^{o(1)}$.  Such a loss is negligible on the
scale of powers of $N$ but need not be bounded.

\begin{lemma}[Endpoint coercivity at $a=\frac12$]
\label{lem:positive-endpoint-coercivity}
There is a sequence $\gamma_N=N^{-o(1)}$ such that
\begin{equation}\label{eq:positive-endpoint-coercivity}
 T_{N,1/2}\ge\gamma_N\id_N.
\end{equation}
\end{lemma}

\begin{proof}
Write $T_{N,1/2}$ in the decomposition
$\C e_0\oplus\operatorname{span}\{e_1,\ldots,e_{N-1}\}$ as
\begin{equation}\label{eq:block-positive-endpoint}
 T_{N,1/2}
 =\begin{pmatrix} a_N&b_N^*\\ b_N&R_N\end{pmatrix},
 \qquad
 R_N=T_{N-1,5/2}.
\end{equation}
The moment representation
\[
 \Hh_{a,N}=\int_0^1 x^{a-1}
 (1,x,\ldots,x^{N-1})^{\mathsf T}
 (1,x,\ldots,x^{N-1})\,dx
\]
shows that $\Hh_{N,5/2}\le\Hh_{N,2}$.  The classical finite Hilbert-matrix
estimate
\begin{equation}\label{eq:Wilf-gap}
 \pi-\|\Hh_{N,2}\|\ge \frac{c}{\log^2(N+1)}
\end{equation}
for all sufficiently large $n$ therefore gives
\begin{equation}\label{eq:R-gap}
 R_N\ge\frac{c}{\log^2(N+1)}\id_{N-1}.
\end{equation}
The sharper asymptotic behind \eqref{eq:Wilf-gap} is due to de Bruijn and
Wilf; only this lower bound is needed here \cite{Wilf}. [TODO: Clarify
  this, "shar per asymptotic vs "only this"?]

Let
\[
 \rho_N:=a_N-b_N^*R_N^{-1}b_N
 =\frac{\det T_{N,1/2}}{\det T_{N-1,5/2}}
\]
be the scalar Schur complement.  Both determinants are positive.  Otte's
endpoint theorem, applied once with parameter $1/2$ and once with parameter
$5/2$, gives
\begin{equation}\label{eq:rho-subpoly}
 \log\rho_N=o(\log N),
 \qquad
 \rho_N^{-1}=N^{o(1)}.
\end{equation}
The second assertion follows from the first because $\rho_N>0$.

Moreover, $\|b_N\|$ is bounded uniformly in $N$.  Put
\[
 U_N:=\begin{pmatrix}1&0\\R_N^{-1}b_N&\id\end{pmatrix}.
\]
The Schur factorization may be written as the congruence
\begin{equation}\label{eq:endpoint-LDU}
 T_{N,1/2}
 =U_N^*
   \begin{pmatrix}\rho_N&0\\0&R_N\end{pmatrix}
  U_N.
\end{equation}
By \eqref{eq:R-gap},
$\|R_N^{-1}b_N\|=O(\log^2(N+1))$, and hence
\begin{equation}\label{eq:endpoint-U-bound}
 \|U_N\|+\|U_N^{-1}\|=O(\log^2(N+1)).
\end{equation}
For every vector $z$,
\[
 \langle z,T_{N,1/2}z\rangle
 \ge
 \min\left\{\rho_N,\frac{c}{\log^2(N+1)}\right\}
 \|U_Nz\|^2.
\]
Using $\|U_Nz\|\ge\|U_N^{-1}\|^{-1}\|z\|$ and
\eqref{eq:rho-subpoly}--\eqref{eq:endpoint-U-bound}, the right-hand side is
bounded below by $\gamma_N\|z\|^2$ for a sequence
$\gamma_N=N^{-o(1)}$.  This proves
\eqref{eq:positive-endpoint-coercivity}.
\end{proof}

The other base parameter, $a=-\frac12$, is not covered by Otte's theorem.  It
can nevertheless be reduced to the two positive parameters $a=\frac12$ and
$a=\frac32$ by elementary eigenvalue interlacing.

\begin{proposition}[The base negative threshold]\label{prop:negative-half-base}
As $N\to\infty$,
\begin{equation}\label{eq:negative-half-det}
 \log\det T_{N,-1/2}=-\frac38\log N+o(\log N).
\end{equation}
Moreover, there is a sequence $\widetilde\gamma_N=N^{-o(1)}$ such that
\begin{equation}\label{eq:negative-half-coercivity}
 T_{N,-1/2}\ge\widetilde\gamma_N\id_N.
\end{equation}
\end{proposition}

\begin{proof}
Set $L_N:=K_{N,-1/2}$.  Its lower-right principal submatrix is
$K_{N-1,3/2}>0$, whereas
$\langle e_0,L_Ne_0\rangle=-2/\pi<0$.  Cauchy interlacing therefore shows that
$L_N$ has exactly one negative eigenvalue.  Write its eigenvalues as
\begin{equation}\label{eq:L-eigenvalues}
 \lambda_1\ge\cdots\ge\lambda_{N-1}>0>\lambda_N.
\end{equation}

The infinite operator $K_{-1/2}$ has the eigenvalue $-1$ with eigenvector given
by the Taylor coefficients of
\begin{equation}\label{eq:negative-half-eigenvector}
 (1-z)^{1/2}.
\end{equation}
These coefficients are $O(n^{-3/2})$.  If $u$ is the normalized coefficient
vector, then $\|Q_Nu\|=O(N^{-1})$, and hence
\[
 \|(L_N+\id_N)P_Nu\|
 =\|P_NK_{-1/2}Q_Nu\|=O(N^{-1}).
\]
Since $\|P_Nu\|\to1$, the spectral theorem implies that the distance from
$-1$ to the spectrum of $L_N$ is $O(N^{-1})$.  The only negative eigenvalue is
$\lambda_N$, while all other eigenvalues are positive.  Hence
\begin{equation}\label{eq:negative-eigenvalue-rate}
 \lambda_N=-1+O(N^{-1}).
\end{equation}

Delete the first row of $L_N$, and denote the resulting
$(N-1)\times N$ matrix by $B_N$.  After reindexing, $B_N$ is the first
$N-1$ rows of $K_{N,1/2}$.  For $j=1,\ldots,N-1$, singular-value
interlacing under deletion of one row and the ideal property give
\begin{equation}\label{eq:singular-interlace-full}
 \sigma_{j+1}(L_N)
 \le \sigma_j(B_N)
 \le \sigma_j(K_{N,1/2}).
\end{equation}
In particular,
\begin{equation}\label{eq:singular-interlace-base}
 \sigma_2(L_N)
 \le\|K_{N,1/2}\|
 \le1-\gamma_N,
\end{equation}
where $\gamma_N=N^{-o(1)}$ is supplied by
Lemma~\ref{lem:positive-endpoint-coercivity}.  Since
$N^{-1}=o(\gamma_N)$, equations \eqref{eq:negative-eigenvalue-rate} and
\eqref{eq:singular-interlace-base} imply that $|\lambda_N|$ is the largest
singular value of $L_N$ for all sufficiently large $N$.

Let $\mu_1\ge\cdots\ge\mu_N>0$ be the eigenvalues of $K_{N,1/2}$.  Once
$|\lambda_N|$ is the largest singular value, the remaining singular values of
$L_N$ are $\lambda_1,\ldots,\lambda_{N-1}$ in that order.  Thus
\eqref{eq:singular-interlace-full} gives the explicit comparison
\begin{equation}\label{eq:positive-eigenvalue-comparison}
 \lambda_j=\sigma_{j+1}(L_N)
 \le\sigma_j(B_N)
 \le\mu_j,
 \qquad j=1,\ldots,N-1.
\end{equation}
The normalized infinite operator $K_{-1/2}$ is a contraction, and so is its
compression $L_N$; in particular, $-1\le\lambda_N<0$.  Consequently,
\begin{align}
 \det(\id_N-L_N)
 &=(1-\lambda_N)\prod_{j=1}^{N-1}(1-\lambda_j)\notag\\
 &\ge\prod_{j=1}^{N-1}(1-\mu_j)
 \ge\det(\id_N-K_{N,1/2}).
 \label{eq:base-lower-bound}
\end{align}
For the reverse estimate, let
$\nu_1\ge\cdots\ge\nu_{N-1}>0$ be the eigenvalues of the lower-right
principal submatrix $K_{N-1,3/2}$.  Cauchy interlacing gives
$\lambda_j\ge\nu_j$ for $j=1,\ldots,N-1$.  Since
$1-\lambda_N\le2$, it follows that
\begin{equation}\label{eq:base-upper-bound}
 \det(\id_N-L_N)
 \le2\prod_{j=1}^{N-1}(1-\nu_j)
 =2\det(\id_{N-1}-K_{N-1,3/2}).
\end{equation}
Otte's endpoint theorem applied to $a=\frac12$ and $a=\frac32$ in
\eqref{eq:base-lower-bound}--\eqref{eq:base-upper-bound} proves
\eqref{eq:negative-half-det}.

Finally, \eqref{eq:positive-eigenvalue-comparison} and
Lemma~\ref{lem:positive-endpoint-coercivity} imply
$\lambda_1\le1-\gamma_N$.  Since $1-\lambda_N>1$, this proves
\eqref{eq:negative-half-coercivity}.
\end{proof}

We next replace the uniform coercivity in
Lemma~\ref{lem:edge-extraction} by the subpolynomial estimates just obtained.
The following formulation is adapted to the Schur complements used below.

\begin{lemma}[Threshold tail extraction]\label{lem:threshold-tail-extraction}
Let $B\ge0$ be a bounded injective operator on $\ell^2(\Nzero)$ and suppose
that
\begin{equation}\label{eq:threshold-B-coercivity}
 P_MB P_M\ge\gamma_M P_M,
 \qquad
 \gamma_M=M^{-o(1)}.
\end{equation}
Let $L:\C^d\to\ell^2(\Nzero)$ be injective, put
$G_M:=L^*Q_ML$, and assume that $G_M$ is positive definite for all sufficiently
large $M$.  Suppose further that there are fixed constants $A>1$, $c>0$, and
$C<\infty$ such that
\begin{equation}\label{eq:scale-separation}
 L^*Q_{\lfloor M^A\rfloor}L
 \le C M^{-c}G_M
\end{equation}
for all sufficiently large $M$.  Define a positive $d\times d$ matrix by
\begin{equation}\label{eq:threshold-S-def}
 x^*S_Mx
 :=\inf_{y\in P_M\ell^2}
 \langle y-Lx,B(y-Lx)\rangle.
\end{equation}
Then
\begin{equation}\label{eq:threshold-S-asymptotic}
 \log\det S_M=\log\det G_M+o(\log M).
\end{equation}
\end{lemma}

\begin{proof}
Choosing $y=P_MLx$ in \eqref{eq:threshold-S-def} gives
\begin{equation}\label{eq:threshold-upper}
 S_M\le\|B\|G_M.
\end{equation}
For the lower bound, fix $y\in P_M\ell^2$, put
$w:=y-Lx$, and let $M_A:=\lfloor M^A\rfloor$.  By the triangle inequality,
\begin{align*}
 \|B^{1/2}w\|
 &\ge \|B^{1/2}P_{M_A}w\|
      -\|B^{1/2}Q_{M_A}w\|\\
 &\ge \gamma_{M_A}^{1/2}\|P_{M_A}w\|
      -\|B\|^{1/2}\|Q_{M_A}Lx\|.
\end{align*}
Because $y$ is supported in the first $M$ coordinates,
\[
 \|P_{M_A}w\|
 \ge\|(P_{M_A}-P_M)Lx\|.
\]
Condition \eqref{eq:scale-separation} therefore implies, uniformly in $x$,
\[
 \|(P_{M_A}-P_M)Lx\|^2
 \ge(1-CM^{-c})x^*G_Mx,
 \qquad
 \|Q_{M_A}Lx\|^2\le CM^{-c}x^*G_Mx.
\]
Since $\gamma_{M_A}=M^{-o(1)}$, one has
$M^{-c/2}\gamma_{M_A}^{-1/2}=M^{-c/2+o(1)}\to0$.  Hence the second term in
the lower bound for $\|B^{1/2}w\|$ is eventually at most one half of the
first, uniformly in $x$.  Taking the infimum over $y$ yields
\begin{equation}\label{eq:threshold-lower}
 S_M\ge\frac14\gamma_{M_A}G_M
\end{equation}
for all sufficiently large $M$.  Since $d$ is fixed,
\eqref{eq:threshold-upper} and \eqref{eq:threshold-lower} imply
\[
 d\log\gamma_{M_A}-d\log4
 \le \log\det S_M-\log\det G_M
 \le d\log\|B\|.
\]
Finally, $\log\gamma_{M_A}=o(\log M_A)=o(\log M)$, proving
\eqref{eq:threshold-S-asymptotic}.
\end{proof}

For the vectors in Lemma~\ref{lem:eigenvector-tails}, condition
\eqref{eq:scale-separation} follows directly from the Cauchy-Gram asymptotic.
Indeed, if the relevant indices are $j_1<\cdots<j_d$ and
\[
 D_M:=\operatorname{diag}
 \bigl(M^{\delta+j_1+1/2},\ldots,M^{\delta+j_d+1/2}\bigr),
\]
then \eqref{eq:tail-pairing} gives
\[
 G_M=D_M(C+o(1))D_M
\]
with $C>0$.  At $\lfloor M^A\rfloor$ the corresponding diagonal matrix is
$D_ME_M(\id+o(1))$, where
\[
 E_M:=\operatorname{diag}
 \bigl(M^{(A-1)(\delta+j_1+1/2)},\ldots,
       M^{(A-1)(\delta+j_d+1/2)}\bigr).
\]
Every exponent in $E_M$ is strictly negative.  Thus there is $c>0$ such
that $\|E_M\|^2\le M^{-c}$.  Since $C$ is positive definite, for all large
$M$ there are constants $0<c_0\le C_0<\infty$ with
\[
 c_0D_M^2\le G_M\le C_0D_M^2
\]
and
\[
 G_{\lfloor M^A\rfloor}
 \le C_0D_ME_M^2D_M
 \le \frac{C_0}{c_0}M^{-c}G_M.
\]
This is \eqref{eq:scale-separation}.  A fixed change of basis conjugates both
Gram matrices by the same invertible matrix and therefore does not affect the
conclusion.

\begin{proposition}[Threshold determinant asymptotics]
\label{prop:threshold-asymptotics}
Let $\delta=-m-\frac12$ with $m=1,2,\ldots$, and set
\begin{equation}\label{eq:threshold-q}
 q:=\frac1\pi\arcsin(\sin(-\pi\delta))=\frac{(-1)^m}{2}.
\end{equation}
Then
\begin{equation}\label{eq:threshold-master}
 \log D_N^\pm(\delta)
 =\left[
 -\frac{q^2\mp q}{2}
 +\sum_{j\in J^\pm_\delta}(2\delta+2j+1)
 \right]\log N+o(\log N).
\end{equation}
Consequently,
\begin{equation}\label{eq:threshold-unfolded}
 \log D_N^\pm(\delta)
 =-\frac{\delta^2\pm\delta}{2}\log N+o(\log N).
\end{equation}
\end{proposition}

\begin{proof}
Put $\alpha=\frac12-m$ and $d:=\lfloor m/2\rfloor$.  Compression to the
coordinates $d,d+1,\ldots$ shifts the generalized Hilbert parameter by $2d$:
\begin{equation}\label{eq:threshold-compression}
 Q_dK_\alpha Q_d\cong K_a,
 \qquad
 a:=\alpha+2d=
 \begin{cases}
  \frac12,&m\ \text{even},\\
  -\frac12,&m\ \text{odd}.
 \end{cases}
\end{equation}
The kernel of $T_\alpha=\id-K_\alpha$ has dimension $d$ and is spanned by the
vectors $u_j$ with
\begin{equation}\label{eq:critical-J}
 0\le j<m,
 \qquad
 j\equiv m\pmod2.
\end{equation}
Indeed, $K_\alpha=(-1)^mA_\alpha$ and
$A_\alpha u_j=(-1)^ju_j$.

Let $E_d:=P_d\ell^2$ and $F_d:=Q_d\ell^2$.  The restriction from
$\Ker T_\alpha$ to $E_d$ is injective: a kernel vector vanishing on $E_d$
would, by \eqref{eq:threshold-compression}, give a vector in $\Ker T_a$, but
$T_{1/2}$ and $T_{-1/2}$ are injective.  Since both spaces have dimension $d$,
the restriction is an isomorphism.  Hence there is a map
$L:E_d\to F_d$ such that
\begin{equation}\label{eq:kernel-graph}
 \Ker T_\alpha=\{x\oplus Lx:x\in E_d\}.
\end{equation}

With respect to $E_d\oplus F_d$, write
\begin{equation}\label{eq:threshold-block-operator}
 T_\alpha=
 \begin{pmatrix}A&C^*\\ C&T_a\end{pmatrix}.
\end{equation}
Since every $x\oplus Lx$ lies in the kernel, the two block equations are
\[
 Cx+T_aLx=0,
 \qquad
 Ax+C^*Lx=0.
\]
Thus $C=-T_aL$ and, using self-adjointness, $A=L^*T_aL$.  Consequently,
\begin{equation}\label{eq:threshold-square-completion}
 \left\langle x\oplus y,T_\alpha(x\oplus y)\right\rangle
 =\langle y-Lx,T_a(y-Lx)\rangle.
\end{equation}

Let $M=N-d$.  The lower-right block of the $N$th finite section is
$T_{M,a}$ and is positive definite.  Taking its Schur complement gives
\begin{equation}\label{eq:threshold-det-factorization}
 \det T_{N,\alpha}=\det T_{M,a}\det S_M,
\end{equation}
where the Schur complement has the equivalent variational description
\begin{equation}\label{eq:threshold-variational}
 x^*S_Mx
 =\inf_{y\in P_MF_d}
 \langle y-Lx,T_a(y-Lx)\rangle.
\end{equation}
Proposition~\ref{prop:negative-half-base} and
Lemma~\ref{lem:positive-endpoint-coercivity} allow us to apply
Lemma~\ref{lem:threshold-tail-extraction} to \eqref{eq:threshold-variational}.
The columns of $L$ are a fixed basis of the kernel in
\eqref{eq:critical-J}; Lemma~\ref{lem:eigenvector-tails} is unchanged by this
fixed change of basis and by deleting finitely many coordinates.  Therefore
\begin{equation}\label{eq:critical-schur-exponent}
 \log\det S_M
 =\left(
   \sum_{\substack{0\le j<m\\j\equiv m\ (2)}}
   (2\delta+2j+1)
  \right)\log M+o(\log M).
\end{equation}%
[TODO: Fix  this last equation. Parens are super large and what does
  the notation mean, modulo?]
Both base determinants in \eqref{eq:threshold-compression} satisfy
\begin{equation}\label{eq:base-three-eighths}
 \log\det T_{M,a}=-\frac38\log M+o(\log M)
\end{equation}
by Otte's theorem when $a=\frac12$ and by
Proposition~\ref{prop:negative-half-base} when $a=-\frac12$.  This proves
\eqref{eq:threshold-master} for the critical determinant, because its
principal-branch term is $-3/8$.

For the other determinant the absolutely continuous part of
$\id+K_\alpha$ is bounded below by the identity.  After adding the projection
onto its finite-dimensional kernel at $-1$, the trace-class comparison and
Lemma~\ref{lem:edge-extraction} apply exactly as in Sections~3--5, now with the
reference operator $\id+\pi^{-1}\Hh_1$.  Proposition~\ref{prop:Hilbert-det}
with coupling $-1$ gives the principal term $1/8$, and the remaining kernel
vectors are those with parity opposite to \eqref{eq:critical-J}.  This proves
\eqref{eq:threshold-master} for the second determinant.

Finally, with $\theta_\delta=-\delta=m+\frac12$, equation \eqref{eq:threshold-q} gives
$q=(-1)^m(\theta_\delta-m)$.  The algebra in Lemma~\ref{lem:unfolding} remains valid at
this endpoint and turns \eqref{eq:threshold-master} into
\eqref{eq:threshold-unfolded}.
\end{proof}

Proposition~\ref{prop:threshold-asymptotics} completes the proof of
Theorem~\ref{thm:main}.  The argument determines the coefficient of $\log N$;
it deliberately allows additional factors of the form $N^{o(1)}$, including
possible powers of $\log N$.
